\subsection{The smallest open case: quartic fields and \texorpdfstring{$n=2$}{n=2}}
\label{ssec:quartic}

By \Cref{thm:closure}(ii) the only form of (P2) on a minimal route to the
conjecture is the one for the split families of the CM fields of degree at
least four, and by \Cref{prop:descent}(i) the Weil classes of any family of
such a field in a higher dimension give those of its families at $n=2$. For a
quartic field these are families of abelian eightfolds, and they are the
smallest families of a field of degree at least four whose Weil classes are
not known to be algebraic: we know of no quartic CM field $F$ and no
discriminant $\delta$ for which $\mathbf{W}(F,2,\delta)$ is proved. This
subsection determines how far the available mechanisms reach there. Divisor
classes reach the Weil classes exactly on a Noether--Lefschetz locus, which
has codimension two when $\delta$ is trivial (\Cref{prop:quarticnl},
\Cref{cor:quarticbase}); for biquadratic $F$ Markman's theorem reaches them on
a further locus of dimension four (\Cref{prop:quarticscalar}); the Casimir
class of the quadratic form on the weight two part $R_{F}$ below, the natural
class of degree four attached to it, is algebraic exactly when the Weil
classes are (\Cref{prop:quarticcasimir}); correspondences with abelian
varieties whose simple factors have dimension at most seven cannot help at a
member with full Hodge group (\Cref{prop:quarticnogo}); and for a general
shape the corrected criterion asks $\dim\Ext^{2}(E,E)=112$ of one complex
(\Cref{thm:quarticp2prime}). None of this is a new case of the Hodge
conjecture, and $\mathbf{W}(F,2,\delta)$ remains open for every quartic $F$
and every $\delta$. Item (LV) of \Cref{sec:verification} checks the
statements in exact arithmetic; its scripts are in
\texttt{code/attack/quartic\_cm}.

Throughout, $F$ is a quartic CM field and $F_{0}$ its real quadratic subfield,
with real places $\tau_{1},\tau_{2}$, and $(V,H)$ is of $(F,2)$-Weil type with
discriminant $\delta$; so $m=n=2$, the members are abelian eightfolds, and
$\cD_{F}$ has dimension $mn^{2}=8$. By \Cref{prop:cmweil}(i),
$H^{1}(A,\CC)=\bigoplus_{\sigma}V_{\sigma}$ over the four embeddings $\sigma$ of
$F$, with $\dim V_{\sigma}^{1,0}=\dim V_{\sigma}^{0,1}=2$. Put
\[
  R_{F}=\textstyle\bigwedge^{2}_{F}V\;\subseteq\;\bigwedge^{2}_{\QQ}V=H^{2}(A,\QQ),
\]
the $\QQ$-subspace with complexification
$\bigoplus_{\sigma}\bigwedge^{2}V_{\sigma}$; for $m=1$ it is the space $R$ of
\Cref{prop:divfails}. It is an $F$-vector space of dimension $6$, of
$\QQ$-dimension $24$, with Hodge numbers $(4,16,4)$. Fix a generator $\Omega$
of the $F$-line $W_{F}=\bigwedge^{4}_{F}V$, so that
$W_{F}=\{a\Omega:a\in F\}$, and define a symmetric $F$-bilinear form $q$ on
$R_{F}$ by $x\wedge_{F}y=q(x,y)\,\Omega$. Let $h$ be the hermitian form on
$R_{F}$ induced by $H$, $h(x_{1}\wedge x_{2},y_{1}\wedge y_{2})
=\det\bigl(H(x_{i},y_{j})\bigr)$, and $h_{4}$ the one induced on
$\bigwedge^{4}_{F}V$; in a basis $e_{1},\dots,e_{4}$ with
$\Omega=e_{1}\wedge e_{2}\wedge e_{3}\wedge e_{4}$, $h_{4}(\Omega,\Omega)$ is
$\det H$, so its class in $F_{0}^{\times}/\Nm_{F/F_{0}}(F^{\times})$ is
$\delta$. With dual bases $(\omega_{k})$, $(\omega_{k}^{*})$ of $F$ for the
trace form, as in \Cref{thm:cmbasepoint}, put
$x\star y=\sum_{k}(\omega_{k}x)\cup(\omega_{k}^{*}y)$ for $x,y\in R_{F}$. Then
\begin{equation}\label{eq:quarticstar}
  x\star y=q(x,y)\,\Omega ,
\end{equation}
because the complexification of $(\omega_{k}x)\cup(\omega_{k}^{*}y)$ is
$\sum_{\sigma,\sigma'}\sigma(\omega_{k})\sigma'(\omega_{k}^{*})\,
x_{\sigma}\wedge y_{\sigma'}$, the sum over $k$ of the coefficient is $1$ for
$\sigma=\sigma'$ and $0$ otherwise, as in the proof of
\Cref{thm:cmbasepoint}(ii), and $\sum_{\sigma}x_{\sigma}\wedge y_{\sigma}$ is
the complexification of $q(x,y)\Omega$. Finally let $\theta_{j}\in H^{1,1}(A)$
be the component of the polarisation that pairs $V_{\sigma}$ with
$V_{\bbar\sigma}$ for the two embeddings $\sigma,\bbar\sigma$ over $\tau_{j}$.
Then $\eta_{t}=\sum_{j}\tau_{j}(t)\,\theta_{j}$ for the classes $\eta_{t}$ of
\Cref{prop:cmweil}(iii), and $\CC[\theta_{1},\theta_{2}]$ is the
complexification of the $\QQ$-algebra generated by the $\eta_{t}$.

\begin{proposition}[Hodge classes, and a quaternion algebra acting on $R_{F}$]
\label{prop:quarticquat}
Let $A$ be a member of an $(F,2,\delta)$-family.
\begin{enumerate}[label=\textup{(\roman*)},leftmargin=2.6em,itemsep=2pt]
\item If $\Hg(A)=G_{F}$, the space of Hodge classes of $A$ has dimension $2$
  in degree two, spanned by the $\eta_{t}$, and $7$ in degree four, spanned by
  the three products of two $\eta_{t}$ and by $W_{F}(A)$, which meets the span
  of those products in $0$.
\item The map $\phi\colon R_{F}\to R_{F}$ defined by $q(x,y)=h(x,\phi y)$
  satisfies $\phi(ax)=\bbar a\,\phi(x)$ for $a\in F$ and
  $\phi^{2}=h_{4}(\Omega,\Omega)^{-1}$, and it commutes with $\SU(V,H)$. So
  $\phi$ and the action of $F$ are Hodge endomorphisms of $R_{F}$ at every
  member, and the Hodge classes in $R_{F}$ form a module over the algebra
  $\cQ_{\delta}=F\oplus F\phi$.
\item $\cQ_{\delta}$ is the cyclic quaternion algebra $(F/F_{0},\delta)$ over
  $F_{0}$, which is split if and only if $\delta$ is trivial. If
  $\Hg(A)=G_{F}$, then $\cQ_{\delta}$ is the whole algebra of Hodge
  endomorphisms of $R_{F}$.
\item If $\delta$ is trivial, then $R_{F}=W\otimes_{F_{0}}F$ for a rational
  sub-Hodge structure $W$, stable under $F_{0}$, of dimension $12$ and with
  Hodge numbers $(2,8,2)$.
\end{enumerate}
\end{proposition}

\begin{proof}
(i) As in the proof of \Cref{prop:cmweil}(iii), over $\CC$ the group $G_{F}$
is $\SL(U_{1})\times\SL(U_{2})$ with $\dim U_{j}=4$, and
$H^{1}(A,\CC)=\bigoplus_{j}(U_{j}\oplus U_{j}^{\vee})$. The invariants of
$\SL(U)$ in $\bigwedge^{a}U\otimes\bigwedge^{b}U^{\vee}$ form a line if $a=b$
or $\{a,b\}=\{0,4\}$, and are zero otherwise. In degree two this leaves the
bidegree $(1,1)$ in one factor: two lines, spanned by $\theta_{1}$ and
$\theta_{2}$. In degree four it leaves $(2,2)$ in one factor or $(1,1)$ in
both, three lines spanned by $\theta_{1}^{2}$, $\theta_{1}\theta_{2}$ and
$\theta_{2}^{2}$, which are nonzero because $\dim U_{j}=4$; and $(4,0)$ or
$(0,4)$ in one factor, four lines spanning $W_{F}\otimes\CC$. They lie in
different summands.

(ii) Choose an $H$-orthogonal basis $e_{1},\dots,e_{4}$ of $V$ and put
$a_{i}=H(e_{i},e_{i})\in F_{0}$. Replacing $\Omega$ by $\lambda\Omega$
replaces $\phi$ by $\bbar\lambda^{-1}\phi$ and $h_{4}(\Omega,\Omega)$ by
$\Nm(\lambda)h_{4}(\Omega,\Omega)$, so we may take
$\Omega=e_{1}\wedge e_{2}\wedge e_{3}\wedge e_{4}$, with
$h_{4}(\Omega,\Omega)=a_{1}a_{2}a_{3}a_{4}$. For $J=\{j<j'\}$ put
$e_{J}=e_{j}\wedge e_{j'}$ and $a_{J}=a_{j}a_{j'}$, and let $J^{c}$ be the
complement, with $e_{J^{c}}\wedge e_{J}=\epsilon_{J}\Omega$. The $e_{J}$ form
an $h$-orthogonal $F$-basis with $h(e_{J},e_{J})=a_{J}$, and
$q(e_{K},e_{J})$ is $\epsilon_{J}$ for $K=J^{c}$ and $0$ otherwise, so
$\phi(e_{J})=\epsilon_{J}a_{J^{c}}^{-1}e_{J^{c}}$. Two-forms commute, so
$\epsilon_{J^{c}}=\epsilon_{J}$ and
$\phi^{2}(e_{J})=(a_{J}a_{J^{c}})^{-1}e_{J}=h_{4}(\Omega,\Omega)^{-1}e_{J}$.
Semilinearity holds because $q$ is linear and $h$ antilinear in the second
variable. An element of $\SU(V,H)$ preserves $h$ and fixes $\Omega$, hence
preserves $q$ and commutes with $\phi$, and it commutes with $F$. The Hodge
group of every member lies in $G_{F}$, because the circle defining the Hodge
structure preserves $E$, commutes with $F$ and acts on each $V_{\sigma}$ with
determinant $z^{2}z^{-2}=1$; so $\phi$ and $F$ commute with it, and they are
Hodge endomorphisms.

(iii) The relations $\phi a=\bbar a\phi$ and $\phi^{2}=h_{4}(\Omega,\Omega)^{-1}$
present $F\oplus F\phi$ as the cyclic algebra
$(F/F_{0},h_{4}(\Omega,\Omega)^{-1})$, which is $(F/F_{0},\delta)$ because a
square in $F_{0}$ is a norm; and a cyclic algebra $(F/F_{0},c)$ is split
exactly when $c$ is a norm from $F$. It has $\QQ$-dimension $8$, since $\phi$
is not $F$-linear. If $\Hg(A)=G_{F}$, the Hodge endomorphisms of $R_{F}$ are
the commutant of $G_{F}$. Over $\CC$,
$R_{F}\otimes\CC=\bigoplus_{j}\bigl(\bigwedge^{2}U_{j}\oplus
\bigwedge^{2}U_{j}^{\vee}\bigr)$, where
$\bigwedge^{2}U_{j}\cong\bigwedge^{2}U_{j}^{\vee}$ is irreducible for
$\SL(U_{j})$ and trivial for the other factor, so the commutant is
$M_{2}(\CC)\times M_{2}(\CC)$, of dimension $8$.

(iv) If $\delta$ is trivial, $h_{4}(\Omega,\Omega)=\Nm(b)$ for some $b\in F$,
and $\psi=b\phi$ satisfies $\psi^{2}=b\bbar b\,\phi^{2}=1$. Its fixed space
$W$ is an $F_{0}$-subspace and a sub-Hodge structure, $\psi$ being a Hodge
endomorphism, and an element $\beta\in F$ with $\bbar\beta=-\beta\ne0$ maps it
isomorphically onto the $(-1)$-eigenspace, since $\psi\beta=-\beta\psi$. So
$R_{F}=W\oplus\beta W=W\otimes_{F_{0}}F$, and $W$ has half the dimension and
half the Hodge numbers of $R_{F}$. Item (LV) checks (i), (ii) and the
dimension in (iii) exactly, for several fields and hermitian forms.
\end{proof}

\begin{proposition}[What divisor classes reach]\label{prop:quarticnl}
Let $A$ be a member of an $(F,2,\delta)$-family.
\begin{enumerate}[label=\textup{(\roman*)},leftmargin=2.6em,itemsep=2pt]
\item The following are equivalent: \textup{(a)} $W_{F}(A)$ lies in the
  subalgebra of $H^{*}(A,\QQ)$ generated by $\NS(A)_{\QQ}$; \textup{(b)}
  $W_{F}(A)$ meets that subalgebra in a nonzero class; \textup{(c)} $R_{F}$
  contains a nonzero rational Hodge class. When they hold, $R_{F}$ contains a
  rational Hodge class $t$ with $q(t,t)\ne0$, and
  $W_{F}(A)=\{(ft)\star t:f\in F\}$, a space of polynomials with rational
  coefficients in divisor classes.
\item For a nonzero $\cQ_{\delta}$-submodule $M\subseteq R_{F}$ let
  $\cD(M)\subseteq\cD_{F}$ be the set of points at which $M$ consists of Hodge
  classes, and let $\mathrm{NL}(R_{F})\subseteq\cD_{F}$ be the set of points at
  which \textup{(c)} holds. Then $\mathrm{NL}(R_{F})=\bigcup_{M}\cD(M)$, a
  countable union. Each nonempty $\cD(M)$ is a closed complex submanifold,
  isomorphic to a product of two bounded symmetric domains of type IV and
  dimension $4-d$, where $d=\dim_{F}M$, and $d=1$ occurs only if $\delta$ is
  trivial. So if $\delta$ is trivial, $\mathrm{NL}(R_{F})$ is the union of the
  nonempty $\cD(M)$ with $d=1$, each of dimension exactly $6$; if $\delta$ is
  not trivial, every $\cD(M)$ has dimension at most $4$.
\item $\mathrm{NL}(R_{F})$ contains the base point $s_{0}$ of
  \Cref{thm:cmbasepoint}, and for every $\delta$ it contains the locus of
  products $B_{1}\times B_{2}$ of two abelian fourfolds of $(F,1)$-Weil type,
  which has dimension $4$. The split locus, of members isogenous to
  $X\otimes_{\cO_{F_{0}}}\cO_{F}$ with $X$ an abelian fourfold with
  $F_{0}\subseteq\End^{0}(X)$, has dimension $6$, lies in the family of
  trivial discriminant, and is contained in $\mathrm{NL}(R_{F})$.
\end{enumerate}
\end{proposition}

\begin{proof}
(i) Clearly (a) implies (b). Let $F^{\times}$ act on $H^{*}(A,\QQ)$ by
$\alpha\mapsto\iota(\alpha)^{*}$, which on
$\bigotimes_{\sigma}\bigwedge^{a_{\sigma}}V_{\sigma}$ is the character
$\prod_{\sigma}\sigma^{a_{\sigma}}$. The isotypic parts for the characters
$\sigma^{2}$ in degree two and $\sigma^{4}$ in degree four are $R_{F}$ and
$W_{F}$. The projections $p_{R}$ and $p_{W}$ onto them are rational, the two
sets of characters being Galois orbits, and they are morphisms of Hodge
structures, being $\QQ$-linear combinations of the $\iota(\alpha)^{*}$. The
characters occurring in $H^{2}$ are the products $\rho\rho'$, and the only way
to write $\sigma^{4}$ as a product of two of them is $\sigma^{2}\cdot\sigma^{2}$,
so $p_{W}(D\cup D')=q(p_{R}D,p_{R}D')\,\Omega$ for $D,D'\in H^{2}(A,\QQ)$. If a
nonzero $w\in W_{F}$ equals $\sum_{i}D_{i}\cup D_{i}'$ with
$D_{i},D_{i}'\in\NS(A)_{\QQ}$, then
$w=p_{W}(w)=\sum_{i}q(p_{R}D_{i},p_{R}D_{i}')\,\Omega$, so some $p_{R}D_{i}$ is
a nonzero rational Hodge class in $R_{F}$, and (b) implies (c). Conversely,
let $t\in R_{F}$ be a nonzero rational Hodge class. The part of type $(1,1)$
of $\bigwedge^{2}V_{\sigma}$ is $V_{\sigma}^{1,0}\wedge V_{\sigma}^{0,1}$, and
$H$ is definite of opposite signs on the two factors by \Cref{prop:cmweil}(i),
so $h$ is negative definite on it at each real place, and $h(t,t)\in F_{0}$ is
totally negative. By \Cref{prop:quarticquat}(ii),
$q(t,\phi t)=h(t,\phi^{2}t)=h_{4}(\Omega,\Omega)^{-1}h(t,t)\ne0$, so one of the
rational Hodge classes $t$, $\phi t$, $t+\phi t$, call it $t'$, has
$q(t',t')\ne0$. Each $\omega_{k}ft'$ and each $\omega_{k}^{*}t'$ is a rational
Hodge class in $H^{2}$, hence a divisor class, and by \eqref{eq:quarticstar}
$(ft')\star t'=f\,q(t',t')\,\Omega$, which runs over $W_{F}$ as $f$ runs over
$F$. So (c) implies (a), with the stated formula.

(ii) The Hodge classes in $R_{F}$ form a $\cQ_{\delta}$-module, so
$\mathrm{NL}(R_{F})$ is the union of the $\cD(M)$, and there are countably
many $M$. Fix a real place $\tau$, let $\sigma$ be an embedding over it, and
put $R_{\tau}=R_{F}\otimes_{F_{0},\tau}\RR$, a complex vector space of
dimension $6$ through $F\otimes_{F_{0},\tau}\RR=\CC$, identified with
$\bigwedge^{2}V_{\sigma}$. Since $\det H$ has sign $(-1)^{n}=1$ at $\tau$, the
map $\phi$ induces an antilinear $\phi_{\tau}$ on $R_{\tau}$ with
$\phi_{\tau}^{2}=\lambda^{2}$, $\lambda>0$. Its eigenspace
$W_{\tau}=\{x:\phi_{\tau}x=\lambda x\}$ is a real form of $R_{\tau}$ and a real
sub-Hodge structure with $W_{\tau}\otimes\CC\cong\bigwedge^{2}V_{\sigma}$, of
Hodge numbers $(1,4,1)$, on which $q=\lambda h$ is real of signature $(2,4)$,
positive on the plane $\Pi_{\tau}$ of real classes of type $(2,0)+(0,2)$ and
negative on those of type $(1,1)$. The group $\SU(V_{\tau},H_{\tau})$ acts on
$W_{\tau}$ preserving $q$, through the isogeny $\SU(2,2)\to\mathrm{SO}(2,4)^{0}$,
and the map from $\cD_{2,2}$ to the space of oriented positive planes in
$W_{\tau}$ that sends a point to $\Pi_{\tau}$ is the equivariant isomorphism of
the symmetric domains of this group onto one component,
$\cD^{\mathrm{IV}}(W_{\tau})$, of dimension $4$. For a quadratic space $P$ of
signature $(2,k)$ write likewise $\cD^{\mathrm{IV}}(P)$, of dimension $k$. Let
$M$ have $F$-dimension $d$. Then $M_{\tau}=M\otimes_{F_{0},\tau}\RR$ is stable
under $\phi$ and $F$, so $M_{\tau}=P_{\tau}\otimes\CC$ with
$P_{\tau}=M_{\tau}\cap W_{\tau}$ of real dimension $d$, and $M$ consists of
Hodge classes exactly when every $P_{\tau}$ is of type $(1,1)$, that is,
orthogonal to $\Pi_{\tau}$. So
$\cD(M)=\prod_{\tau}\cD^{\mathrm{IV}}(P_{\tau}^{\perp})$, which is empty unless
every $P_{\tau}$ is negative definite and then has dimension $2(4-d)$. If
$d=1$, then $M=Fx$ with $\phi x=bx$ for some $b\in F$, and
$\phi^{2}x=\bbar bb\,x$ gives $h_{4}(\Omega,\Omega)^{-1}=\Nm(b)$, so $\delta$ is
trivial. If $\delta$ is trivial, $\cQ_{\delta}\cong M_{2}(F_{0})$, every nonzero
module contains one with $d=1$, and $\cD(M)\subseteq\cD(M')$ for
$M'\subseteq M$.

(iii) The classes $\delta_{i}(f)$ of \Cref{thm:cmbasepoint}(ii) have
complexification in $\bigoplus_{\sigma}\bigwedge^{2}V_{\sigma}$, so they lie in
$R_{F}$, and they are Hodge classes at $s_{0}$. For a product,
$V=V_{1}\oplus V_{2}$ with $H$-orthogonal $F$-planes of signature $(1,1)$ at
each place, and $\bigwedge^{2}_{F}V_{1}\subseteq R_{F}$ is $W_{F}(B_{1})$, of
type $(1,1)$ by \Cref{prop:cmweil}(i). Every $\delta$ is the product of the
discriminants of two such planes, for instance of $\mathrm{diag}(1,-\delta)$
and $\mathrm{diag}(1,-1)$, and each factor varies in a family of dimension
$m=2$. On the split locus $H^{1}(A,\QQ)=U\otimes_{F_{0}}F$, with
$U=H^{1}(X,\QQ)$ and the complex structure of $X$ on the first factor. The
polarisation class $e$ of $X$ lies in $\bigwedge^{2}_{F_{0}}U$, the hermitian
form of $A$ is $c$ times the sesquilinear extension of the $F_{0}$-bilinear
form $e$, with $c\in F$ totally imaginary, and so $\det H=c^{4}\,\mathrm{Pf}(e)^{2}$
is a square in $F_{0}$ and $\delta$ is trivial. The class
$e\otimes1\in(\bigwedge^{2}_{F_{0}}U)\otimes_{F_{0}}F=R_{F}$ is fixed by
$\Hg(A)=\Hg(X)\subseteq\Res_{F_{0}/\QQ}\Sp(U,e)$, so it is a nonzero Hodge
class. The polarised abelian fourfolds with real multiplication by $F_{0}$
form families of dimension $[F_{0}:\QQ]\cdot3=6$. Item (LV) checks the split
locus for three quartic fields.
\end{proof}

\begin{corollary}[A base locus of codimension two]\label{cor:quarticbase}
In every $(F,2,\delta)$-family, $\Sigma_{F}\supseteq\mathrm{NL}(R_{F})$, and at
every point of $\mathrm{NL}(R_{F})$ the Weil classes are the polynomials in
divisor classes of \Cref{prop:quarticnl}(i). If $\delta$ is trivial,
$\mathrm{NL}(R_{F})$ is a countable union of closed complex submanifolds of
$\cD_{F}$ of codimension exactly two, and it contains the split locus; if
$\delta$ is not trivial, they have codimension at least four, and the product
loci have codimension four. In both cases $\mathrm{NL}(R_{F})$ contains the
base point $s_{0}$ of \Cref{thm:cmbasepoint} and is stable under
$G_{F}(\QQ)$. At a point outside $\mathrm{NL}(R_{F})$ no nonzero Weil class is
a polynomial in divisor classes.
\end{corollary}

\begin{proof}
Everything except the stability is \Cref{prop:quarticnl}. An element
$g\in G_{F}(\QQ)$ commutes with $F$ and with $\phi$, so it carries
$\cQ_{\delta}$-submodules to $\cQ_{\delta}$-submodules, and
$g\,\cD(M)=\cD(gM)$.
\end{proof}

\Cref{thm:cmbasepoint} supplied one CM point, whose $G_{F}(\QQ)$-orbit is dense
but countable. When $\delta$ is trivial, \Cref{cor:quarticbase} replaces it by
a locus of codimension two on which the Weil classes are balanced products of
divisor classes, as they are at $s_{0}$; nothing is claimed about the rest of
the Hodge ring at its points. It does not decide $\mathbf{W}(F,2,\delta)$:
$\mathrm{NL}(R_{F})$ is a countable union of proper closed submanifolds, hence
meagre, and \Cref{thm:cmpropagation}(iii) allows $\Sigma_{F}$ to be meagre.
The last sentence of the corollary is the analogue of \Cref{prop:divfails}:
no argument through divisor classes reaches a member outside
$\mathrm{NL}(R_{F})$. The codimension is two rather than
$h^{2,0}(R_{F})=4$ (compare \Cref{rem:nlcodim}) because Hodge classes in
$R_{F}$ come in $\cQ_{\delta}$-modules, and a module of $F$-dimension one costs
one condition at each real place.

\begin{figure}[tbp]
\centering
  \includegraphics[width=0.94\textwidth]{fig_quarticloci}
\caption{Where the Weil classes of an $(F,2,\delta)$-family are known to be
algebraic, $F$ a quartic CM field; each panel is the period domain
$\cD_{F}$, of dimension $8$. The grass regions are components of the
Noether--Lefschetz locus $\mathrm{NL}(R_{F})$, on which $W_{F}$ consists of
polynomials in divisor classes (\Cref{prop:quarticnl},
\Cref{cor:quarticbase}): of dimension exactly $6$ when $\delta$ is trivial,
where the quaternion algebra $\cQ_{\delta}$ of \Cref{prop:quarticquat} is
split, and at most $4$ otherwise. The ochre curve is the locus of scalar
extensions of Weil fourfolds, of dimension $4$ when $F$ is biquadratic, on
which Markman's theorem applies (\Cref{prop:quarticscalar}); it is not
contained in $\mathrm{NL}(R_{F})$ but meets it (ring) where the fourfold is a
product of Weil surfaces. The indigo points are CM points; the base point
$s_{0}$ of \Cref{thm:cmbasepoint} lies on $\mathrm{NL}(R_{F})$. The rest of
the domain, the Hodge-generic member included, is open.}
\label{fig:quarticloci}
\end{figure}

\begin{proposition}[Scalar extension from an imaginary quadratic subfield]
\label{prop:quarticscalar}
Let $F=KF_{0}$ with $K$ imaginary quadratic; these are the biquadratic
quartic CM fields, and the only ones containing an imaginary quadratic field.
Let $C$ be an abelian fourfold of $(K,2,\delta)$-Weil type and
$B_{F}=C\otimes_{\cO_{K}}\cO_{F}$, which is of $(F,2,\iota(\delta))$-Weil type by
\Cref{prop:descent}(ii); write $(V_{C},H_{C})$ for the hermitian space of $C$.
For $f\in F$ let $g_{f}\in\Hom(B_{F},C)\otimes\QQ$ be the element inducing
$v\mapsto v\otimes f$ from $H^{1}(C,\QQ)$ to
$H^{1}(B_{F},\QQ)=H^{1}(C,\QQ)\otimes_{K}F$.
\begin{enumerate}[label=\textup{(\roman*)},leftmargin=2.6em,itemsep=2pt]
\item $W_{F}(B_{F})$ lies in the $\QQ$-span of the classes $g_{f}^{*}w$, $f\in F$,
  $w\in W_{K}(C)$, which has dimension $10$. Hence $W_{F}(B_{F})$ is
  algebraic, by Markman's theorem $\mathbf{W}(K,2,\delta)$ \cite{Mar25a}
  (\Cref{prop:separate}(ii)).
\item The members $B_{F}$ form a locus of dimension $4$ in the
  $(F,2,\iota(\delta))$-family. At those with $\Hg(C)=\SU(V_{C},H_{C})$ the
  space $R_{F}$ contains no nonzero Hodge class, so
  this locus is not contained in $\mathrm{NL}(R_{F})$; it meets
  $\mathrm{NL}(R_{F})$, for instance where $C=C_{1}\times C_{2}$ with $C_{1}$,
  $C_{2}$ abelian surfaces of Weil type, and then $B_{F}$ is a product of two
  abelian fourfolds of $(F,1)$-Weil type.
\end{enumerate}
\end{proposition}

\begin{proof}
(i) The map $v\mapsto v\otimes f$ commutes with the complex structures, that of
$B_{F}$ being the one of $C$ tensored with $F$, so $g_{f}$ exists. Over $\CC$,
for each embedding $\kappa$ of $K$ let $V_{\kappa}$ be the $\kappa$-eigenspace
of $K$ on $H^{1}(C,\CC)$ and $U_{\kappa}=F\otimes_{K,\kappa}\CC\cong\CC^{2}$, the
sum of the lines $L_{\sigma}$, $\sigma|\kappa$, on which $F$ acts through
$\sigma$. The $\kappa$-part of $H^{1}(B_{F},\CC)$ is $V_{\kappa}\otimes
U_{\kappa}$ and its $\sigma$-part is $V_{\kappa}\otimes L_{\sigma}$, so the
$\sigma$-component $\bigwedge^{4}V_{\kappa}\otimes L_{\sigma}^{\otimes4}$ of
$W_{F}(B_{F})\otimes\CC$ lies in $\bigwedge^{4}V_{\kappa}\otimes
\mathrm{Sym}^{4}U_{\kappa}\subseteq\bigwedge^{4}(V_{\kappa}\otimes U_{\kappa})$.
The map $g_{f}^{*}$ sends a generator $\alpha_{\kappa}$ of $\bigwedge^{4}V_{\kappa}$
to $\alpha_{\kappa}\otimes f_{\kappa}^{\otimes4}$, with $f_{\kappa}$ the image
of $f$ in $U_{\kappa}$. The fourth powers of the elements of the Zariski
dense subset $F$ of $U_{\kappa}$ span $\mathrm{Sym}^{4}U_{\kappa}$, so the
complexified span of the $g_{f}^{*}w$ is $\bigoplus_{\kappa}\bigwedge^{4}V_{\kappa}
\otimes\mathrm{Sym}^{4}U_{\kappa}$, of dimension $2\cdot5$, and it contains
$W_{F}(B_{F})\otimes\CC$. A rational subspace whose complexification lies in
that of a rational subspace lies in it, and pull-backs of algebraic classes
are algebraic.

(ii) The map $C\mapsto B_{F}$ embeds the four-dimensional domain of the
$K$-family into $\cD_{F}$. As rational Hodge structures
$H^{1}(B_{F},\QQ)\cong H^{1}(C,\QQ)^{2}$, so $\Hg(B_{F})=\Hg(C)$, and
$R_{F}\otimes\CC=\bigoplus_{\sigma}\bigwedge^{2}V_{\kappa}\otimes
L_{\sigma}^{\otimes2}$, with $\bigwedge^{2}V_{\kappa}$ irreducible and
nontrivial for $\SL(V_{\kappa})$; so when $\Hg(C)=\SU(V_{C},H_{C})$, $R_{F}$
has no invariants. For $C=C_{1}\times C_{2}$ one has
$B_{F}=(C_{1}\otimes_{\cO_{K}}\cO_{F})\times(C_{2}\otimes_{\cO_{K}}\cO_{F})$, and
\Cref{prop:quarticnl}(iii) applies. Item (LV) checks (i), and that $R_{F}$
meets the invariants of $\SU(V_{C},H_{C})$ in $0$, exactly for three
biquadratic fields and four hermitian forms.
\end{proof}

This is not a new case of the Hodge conjecture. $B_{F}$ is isogenous to
$C\times C$, the classes are Hodge classes there, and when
$\Hg(C)=\SU(V_{C},H_{C})$ \Cref{prop:powers} already gives the conjecture for
$C\times C$ from \cite{Mar25a}.

\begin{proposition}[The Casimir class]\label{prop:quarticcasimir}
Let $A$ be a member of an $(F,2,\delta)$-family and $c\in F^{\times}$, and let
\[
  C_{q}\in R_{F}\otimes_{\QQ}R_{F}\subseteq H^{2}(A,\QQ)\otimes H^{2}(A,\QQ)
  \subseteq H^{4}(A\times A,\QQ)
\]
be the Casimir element $\sum_{j}u_{j}\otimes u_{j}^{*}$ of the nondegenerate
symmetric form $\Tr_{F/\QQ}(c^{-1}q)$ on $R_{F}$, for a basis $(u_{j})$ and its
dual basis $(u_{j}^{*})$. Then $C_{q}$ is a Hodge class,
$\Delta^{*}C_{q}=6c\,\Omega$ for the diagonal $\Delta\colon A\to A\times A$,
and $C_{q}$ lies in the $\QQ$-span of the classes $f^{*}w$, $w\in W_{F}(A)$,
$f(x,y)=\iota(a)x+\iota(b)y$ with $a,b\in\cO_{F}$. Hence $W_{F}(A)$ is
algebraic if and only if $C_{q}$ is.
\end{proposition}

\begin{proof}
For an $F$-basis $(x_{I})$ of $R_{F}$ and its $q$-dual basis $(x_{I}^{q})$, the
elements $\omega_{k}x_{I}$ and $c\,\omega_{k}^{*}x_{I}^{q}$ form dual bases for
$\Tr_{F/\QQ}(c^{-1}q)$. So by \eqref{eq:quarticstar}
$\Delta^{*}C_{q}=\sum_{I}x_{I}\star(c\,x_{I}^{q})=\sum_{I}c\,\Omega=6c\,\Omega$,
and if $C_{q}$ is algebraic so is $W_{F}(A)$, by \Cref{prop:cmweil}(ii).
Conversely, $\SU(V,H)$ preserves $q$, so $C_{q}$ is invariant under the
diagonal action of $G_{F}$, which contains $\Hg(A\times A)$, and $C_{q}$ is a
Hodge class. Since $q$ pairs $\bigwedge^{2}V_{\sigma}$ only with itself, the
complexification of $C_{q}$ lies in $\bigoplus_{\sigma}\bigwedge^{2}V_{\sigma}
\otimes\bigwedge^{2}V_{\sigma}$, inside the part
$\bigoplus_{\sigma}\bigwedge^{4}(V_{\sigma}\otimes\CC^{2})$ of
$H^{4}(A\times A,\CC)$ on which $F^{\times}$, acting diagonally, has the
characters $\sigma^{4}$. The group $G_{F}$ acts on $V_{\sigma}$ through
$\SL(V_{\sigma})$, whose invariants in $\bigwedge^{4}(V_{\sigma}\otimes\CC^{2})$
are $\bigwedge^{4}V_{\sigma}\otimes\mathrm{Sym}^{4}\CC^{2}$ by Cauchy's
formula, so the $G_{F}$-invariants in that part form the rational space
$W_{F}(A)\otimes_{F}\mathrm{Sym}^{4}_{F}(F^{2})$, of dimension $20$. The map
$f^{*}$ sends $w$ to $w\otimes(a,b)^{\otimes4}$, and the fourth powers of the
elements of $\cO_{F}^{2}$ span $\mathrm{Sym}^{4}_{F}(F^{2})$; so the classes
$f^{*}w$ span that space, and $C_{q}$ is algebraic if $W_{F}(A)$ is. Item (LV)
checks the three identities exactly for four pairs of field and hermitian
form.
\end{proof}

\begin{remark}[What the exceptional isogeny gives]\label{rem:quarticisogeny}
By the proof of \Cref{prop:quarticnl}(ii), at each real place $R_{F}$ becomes
$W_{\tau}\otimes\CC$ with $W_{\tau}$ of K3 type, $\SU(V_{\tau},H_{\tau})$ acting
through $\SU(2,2)\to\mathrm{SO}(2,4)$, and when $\delta$ is trivial the
$W_{\tau}$ are the real components of the rational structure $W$ of
\Cref{prop:quarticquat}(iv). A route of Kuga--Satake type would start from
this weight two structure and has to return to degree four, on $A$ or on
$A\times A$; \Cref{prop:quarticcasimir} says that the natural class it meets
there, the Casimir class of $q$, is algebraic exactly when the Weil classes
are. So this mechanism restates $\mathbf{W}(F,2,\delta)$ at each member and
does not reduce it. One can go one step further, which we only sketch. Under
the identification $R_{F}\otimes R_{F}\cong\End(R_{F})$ given by the
intersection pairing, $C_{q}$ corresponds to $\phi$ up to multiplication by an
element of $F$, provided that pairing restricts to $R_{F}$ as
$\Tr_{F/\QQ}(a\,h)$ for some $a\in F_{0}$; since the K\"unneth components and
the inverse of the Lefschetz operator of an abelian variety are algebraic
\cite[\S2]{Kle68}, $C_{q}$ would then be algebraic exactly when $\phi$ is
induced by an algebraic correspondence. We have not checked the form of the
pairing, and nothing in this paper uses the sketch.
\end{remark}

\begin{proposition}[No correspondence from a small abelian variety]
\label{prop:quarticnogo}
Let $A$ be a member with $\Hg(A)=G_{F}$, and let $X$ be an abelian variety
such that $\Hg(X)$ has no normal subgroup isogenous to $G_{F}$; by the proof,
this holds when every simple factor of $X$ has dimension at most $7$, in
particular when $\dim X\le7$, and it holds when $X$ is of CM type. Then
$\Hg(A\times X)=\Hg(A)\times\Hg(X)$, and
$\Hdg^{*}(A^{a}\times X^{b})=\Hdg^{*}(A^{a})\otimes\Hdg^{*}(X^{b})$ for all
$a,b\ge0$. So the class of an algebraic cycle $Z$ on $A^{a}\times X^{b}$ is
$\sum_{i}\beta_{i}\otimes\xi_{i}$ with $\beta_{i}\in\Hdg^{*}(A^{a})$ and
linearly independent $\xi_{i}\in\Hdg^{*}(X^{b})$; the correspondence $Z$
carries $H^{*}(X^{b},\QQ)$ into the span of the $\beta_{i}$; and if the Hodge
conjecture holds for $X^{b}$, every $\beta_{i}$ is algebraic.
\end{proposition}

\begin{proof}
The group $G_{F}=\Res_{F_{0}/\QQ}\SU(V,H)$ is almost simple over $\QQ$. By
Goursat's lemma the kernel $N$ of the projection $\Hg(A\times X)\to\Hg(X)$
maps onto a normal subgroup of $\Hg(A)=G_{F}$. If that subgroup is $G_{F}$,
then $\Hg(A\times X)$ contains $G_{F}\times1$ and is the product. Otherwise
$N$ is finite, $\Hg(X)$ has a normal subgroup $H$ isogenous to $G_{F}$, and
the cocharacter $\mu$ defining the Hodge structure of $X$ projects through
$H\to G_{F}$ to that of $A$. We show that this is impossible under the
hypothesis on the simple factors. Over $\CC$, $H$ is $\SL_{4}\times\SL_{4}$,
and $\mu$ projects in each factor to
$\mathrm{diag}(\frac12,\frac12,-\frac12,-\frac12)$ modulo the centre. Since
$\mu$ acts on $H^{1}(X,\CC)$ with two weights differing by one, the weights of
its projection spread by at most one on every irreducible constituent for
$H$; for the representation of $\SL_{4}$ of highest weight
$a\varpi_{1}+b\varpi_{2}+c\varpi_{3}$ the spread is $a+2b+c$, and spreads add
over the two factors. So every nontrivial constituent is the standard
representation of one factor or its dual. These four are permuted
transitively by $\Gal(\bbar\QQ/\QQ)$, as the embeddings of $F$ are; so a simple
factor $Y$ of $X$ on whose $H^{1}$ the group $H$ acts nontrivially has
$\dim_{\QQ}H^{1}(Y,\QQ)\ge16$, that is $\dim Y\ge8$. A torus has no such $H$.
Next, $\Hg(A^{a}\times X^{b})$ is $\Hg(A)\times\Hg(X)$ acting factorwise on
$H^{*}(A^{a})\otimes H^{*}(X^{b})$, and the invariants of a product of groups
in a tensor product of representations of the factors are the tensor product
of the invariants. Finally, the intersection pairing is perfect on
$\Hdg^{*}(X^{b})$, the invariants being orthogonal to the other isotypic
parts, so there are $\xi_{j}^{*}\in\Hdg^{*}(X^{b})$ with
$\int\xi_{i}\xi_{j}^{*}=\delta_{ij}$; they are algebraic when the conjecture
holds for $X^{b}$, and then $\beta_{j}=pr_{A^{a}*}(Z\cdot pr_{X^{b}}^{*}\xi_{j}^{*})$
is algebraic.
\end{proof}

So a correspondence with such an $X$ can produce a Weil class of $A$ only from
an algebraic cycle on some $A^{a}\times X^{b}$ that already yields one on
$A^{a}$. At a member with full Hodge group this excludes every abelian
variety of dimension at most seven, among them the fourfolds and sixfolds of
Weil type of the imaginary quadratic fields on which Markman's theorems are
proved, and every CM abelian variety. The same argument, with real points in
place of dimensions, excludes rational Hodge structures $Y$ of K3 type: by
Zarhin's theorem \cite{Zar83} the Hodge group of such a $Y$ is the restriction
of scalars of an orthogonal or unitary group over the endomorphism field of
its transcendental part, noncompact at exactly one real place, so its real Lie
algebra has at most one simple factor isomorphic to $\mathfrak{su}(2,2)$,
while that of $G_{F}$ has two; hence $\Hg(A\times Y)=\Hg(A)\times\Hg(Y)$.
This part rests on Zarhin's theorem, which we quote.

\begin{theorem}[The corrected criterion for a quartic field at $n=2$]
\label{thm:quarticp2prime}
Let $A$ be a member of an $(F,2,\delta)$-family, let
$\omega=\sum_{\sigma}u_{\sigma}\alpha_{\sigma}$, with $\alpha_{\sigma}$ spanning
$\bigwedge^{4}V_{\sigma}$ and every $u_{\sigma}\ne0$, as holds for every
nonzero rational Weil class, let $p\in\CC[\theta_{1},\theta_{2}]$, put
$\gamma=\omega+p$, and let $r(\gamma)$ be the rank of the contraction
$HT^{2}(A)\to H^{*}(A,\CC)$, $\xi\mapsto\xi\lrcorner\gamma$. Write
$T_{\sigma}=(V_{\sigma}^{1,0})^{\vee}\subseteq T_{0}A$.
\begin{enumerate}[label=\textup{(\roman*)},leftmargin=2.6em,itemsep=2pt]
\item $\dim HT^{2}(A)=28+64+28=120$, and
  \[
    \Ann_{HT^{2}}(\omega)=\textstyle\bigoplus_{\sigma}
    \Hom\bigl(V_{\sigma}^{1,0},V_{\sigma}^{0,1}\bigr)\;\oplus\;
    \bigoplus_{\{\sigma,\sigma'\},\,\sigma\ne\sigma'}T_{\sigma}\otimes T_{\sigma'},
  \]
  of dimension $0+16+24$: the $F$-linear first order deformations of $A$ and
  the bivectors mixing two embeddings. So $r(\omega)=80$. Unlike the case
  $m=1$ of \Cref{thm:p2forces}(i), no nonzero element of $H^{0,2}(A)$ kills
  $\omega$.
\item $r(\gamma)\le112$ for every $p$. For a general $p$, $r(\gamma)=112$, and
  whenever $r(\gamma)=112$ the annihilator of $\gamma$ is the tangent space
  $T_{[A]}\cD_{F}$ of the family, of dimension $8$, inside $H^{1}(T_{A})$. The
  value $112$ is also taken at a general element, with nonzero component in
  $W_{F}\otimes\CC$, of the complexified Hodge ring of a member with
  $\Hg(A)=G_{F}$.
\item $r(\gamma)\ge68$ for every $p$.
\item At the classes $\omega$ and coefficients computed in item (LV): for $p$
  a multiple of one monomial $\theta_{1}^{i}\theta_{2}^{j}$, $0\le i,j\le4$,
  the rank $r(\gamma)$ takes the six values $80$, $84$, $88$, $104$, $108$,
  $112$; for $p$ a sum of two such terms the values $92$ and $96$ occur as
  well, for instance $r(\omega+\theta_{2}+t\theta_{1}^{2})=92$ and
  $r(\omega+\theta_{1}^{2}+t\theta_{2}^{2})=96$ for each of the seven nonzero
  values of $t$ tested. No value below $80$ has been found, and whether the
  bound of \textup{(iii)} is attained is open.
\item Let $E$ be a perfect complex on $A$ with $\ch(E)=N\omega_{1}+P$, where
  $N\ne0$, $\omega_{1}$ is a nonzero rational Weil class and $P$ is a
  polynomial with rational coefficients in the classes $\eta_{t}$. Then
  $\dim\Ext^{2}(E,E)\ge r(\ch E)\ge68$. If $\dim\Ext^{2}(E,E)=r(\ch E)$, then
  $\operatorname{ev}_{E}$ is onto with kernel $\Ann_{HT^{2}}(\ch E)$ and
  $\sigma_{E}$ is injective, so that $\mathbf{W}(F,2,\delta)$ follows by
  \Cref{thm:cmpropagation}(v).
\end{enumerate}
\end{theorem}

\begin{proof}
(i) $HT^{2}(A)=\bigwedge^{2}H^{0,1}\oplus H^{0,1}\otimes T_{0}A\oplus
\bigwedge^{2}T_{0}A$, with $H^{0,1}$ and $T_{0}A$ of dimension $8$. As in the
proof of \Cref{thm:p2forces}(i) the three summands map into different Hodge
components, so the annihilator is the sum of the three annihilators; and the
torus $(\CC^{\times})^{4}$ acting on each $V_{\sigma}$ by scalars separates the
images further. A block $V_{\rho}^{0,1}\wedge V_{\rho'}^{0,1}$ of
$\bigwedge^{2}H^{0,1}$ kills $\alpha_{\sigma}$ exactly when
$\sigma\in\{\rho,\rho'\}$, and otherwise maps injectively into a weight space
that no other pair of a block and an embedding reaches; there are four
embeddings, so some $\sigma$ lies outside $\{\rho,\rho'\}$, and no nonzero
element of $\bigwedge^{2}H^{0,1}$ kills $\omega$. A block
$\Hom(V_{\rho}^{1,0},V_{\rho'}^{0,1})$ of $H^{1}(T_{A})$ kills $\alpha_{\sigma}$
for $\sigma\ne\rho$; it kills $\alpha_{\rho}$ when $\rho'=\rho$, because
$V_{\rho}^{0,1}$ is already exhausted in $\alpha_{\rho}$, and maps it
injectively into its own weight space when $\rho'\ne\rho$. A bivector in
$T_{\rho}\otimes T_{\rho'}$ with $\rho\ne\rho'$ kills every $\alpha_{\sigma}$,
whose two factors of type $(1,0)$ both lie in $V_{\sigma}^{1,0}$, and
$\bigwedge^{2}T_{\sigma}$ maps $\alpha_{\sigma}$ to a nonzero multiple of a
generator of $\bigwedge^{2}V_{\sigma}^{0,1}$. The dimensions are $0$,
$4\cdot4$ and $6\cdot4$.

(ii) The classes $\theta_{j}$ and the rational Weil classes stay of Hodge type
along $\cD_{F}$, so each is killed by $T_{[A]}\cD_{F}\subseteq H^{1}(T_{A})$, as
in \Cref{prop:p2primelocus}(i); hence $r(\gamma)\le120-8$, and when equality
holds the annihilator, of dimension $8$, is $T_{[A]}\cD_{F}$. Take a basis of
$H^{1}(A,\CC)$ adapted to the $V_{\sigma}$ and to the Hodge decomposition, in
which the polarisation pairs the basis of $V_{\sigma}^{1,0}$ with that of
$V_{\bbar\sigma}^{0,1}$. The torus acting diagonally on it contains a subtorus
that multiplies each $\theta_{j}$ by a scalar and each $\alpha_{\sigma}$ by an
arbitrary scalar, and contraction is equivariant; so $r(\omega+p)$ depends on
$\omega$ only through a rescaling of the variables of $p$, and each statement
of the theorem holds for every $\omega$ once it holds for one. The rank is
lower semicontinuous in $p$, and item (LV) finds $r(\gamma)=112$ exactly at
explicit rational $\omega$ and $p$, and at explicit elements of the Hodge ring
with nonzero Weil component.

(iii) In such a basis each basis element of $HT^{2}$ carries a monomial to a
multiple of a monomial. Let $G$ be the set of monomials that are not of the
form $\xi_{0}\lrcorner\mu$ for a basis element $\xi_{0}$ and a monomial $\mu$
occurring in some $\theta_{1}^{i}\theta_{2}^{j}$. The coordinates of
$\xi\lrcorner\gamma$ along $G$ are those of $\xi\lrcorner\omega$, whatever $p$
is, and the map $\xi\mapsto(\xi\lrcorner\omega)|_{G}$ has rank $68$, computed
exactly in item (LV); the rank of a projection bounds $r(\gamma)$ from below.
The full torus acts on the monomials of $G$ by scalars and moves any $\omega$
to any other, so the rank $68$ does not depend on $\omega$.

(iv) is computed exactly in item (LV), on the $25$ monomials and on all $300$
pairs of them with seven ratios of coefficients. By the torus of (ii) the
value on a monomial of total degree other than two does not depend on its
coefficient or on $\omega$; for $\theta_{1}^{2}$, $\theta_{1}\theta_{2}$ and
$\theta_{2}^{2}$ the ratio of the coefficient to $\omega$ is an invariant, as
at the exceptional ratio of \Cref{rem:p2primeexceptional}, and only the ratios
computed are asserted.

(v) By the commutative square $\sigma_{E}\circ\operatorname{ev}_{E}=(\,\cdot\,)
\lrcorner\ch(E)$ of \Cref{setup:criterion}, which holds on every abelian
variety, $\dim\Ext^{2}(E,E)\ge\operatorname{rk}\operatorname{ev}_{E}\ge
r(\ch E)$, and equality forces $\operatorname{ev}_{E}$ to be onto with kernel
the annihilator and $\sigma_{E}$ to be injective, as in
\Cref{thm:p2primenumber}(iv); and $r(\ch E)\ge68$ by (iii), since the
complexification of $\ch(E)$ is a $\gamma$ as above with $\omega=N\omega_{1}$.
The last implication is \Cref{thm:cmpropagation}(v), whose character is
exactly of this form.
\end{proof}

At $m=1$ and $n=2$ the corresponding numbers are $28$, $12$ and $24$
(\Cref{thm:p2primenumber}); in both cases the general value is the dimension
of $HT^{2}$ minus that of the family. For the pure shape the argument of
\Cref{thm:p2forces}(ii) and (iv) transfers with the annihilator of (i): an
injective $\sigma_{E}$ would force $A_{\theta}A_{\theta'}=0$ for
$\theta\in T_{\rho}$, $\theta'\in T_{\rho'}$, $\rho\ne\rho'$, and a smooth
component of the support of codimension at least two along which $E$ is
locally free of positive rank would be tangent to three of the four
$T_{\rho}$. We have not extended \Cref{thm:p2false} to fields of degree four,
so we do not know whether the pure shape is excluded here as it is for $m=1$;
nothing below uses it.

\begin{remark}[What is missing at $n=2$]\label{rem:quarticgap}
The Weil classes of an $(F,2,\delta)$-family are known to be algebraic at the
CM points, on $\mathrm{NL}(R_{F})$, and, for biquadratic $F$, on the
scalar-extension loci, in a period domain of dimension $8$. Each mechanism
examined here either stops on a locus of codimension at least two, or is
equivalent to the statement to be proved, or is excluded at a member with
full Hodge group. By \Cref{prop:descent}(i), $\mathbf{W}(F,3,\delta')$ for a
single discriminant $\delta'$, the trivial one for instance, would give
$\mathbf{W}(F,2,\delta)$ for every $\delta$; by \Cref{prop:descent}(ii) the
families of $F$ at $n=2$ give nothing new for the imaginary quadratic subfields
of $F$, if it has any, whose fourfolds are settled by Markman. What is missing
is propagation, in the form of \Cref{thm:quarticp2prime}(v): one perfect
complex on one member, with Chern character a nonzero multiple of a rational
Weil class plus a polynomial in the $\eta_{t}$, and with
$\dim\Ext^{2}(E,E)=r(\ch E)$, which is $112$ for a general shape; every such
complex has $\dim\Ext^{2}(E,E)\ge68$. When $\delta$ is trivial the secant
sheaves of \cite{Mar25c}, on $X\times\widehat{X}$ with $X$ an abelian fourfold
with real multiplication by $F_{0}$, are natural candidates; we do not know
whether their Chern characters have this form, or whether they are
semiregular. Even $\mathbf{W}(F,2,\delta)$ for every quartic $F$ and every
$\delta$ would be only the first case of the conclusion of
$\mathrm{P2}_{\mathrm{split}}$, which concerns every field of degree at least
four in infinitely many dimensions, and it would leave (F2) and (F3$'$) of
\Cref{cor:fourremain} untouched.
\end{remark}
